\chapter{State Of The Art} \label{chap:ch3}

In this chapter, we analyse the current state of the art described
by the documents gathered in the literature review (Chapter~\ref{chap:ch2}).
We will review different approaches and techniques applied in Edge and \ac{IoT}
environments to achieve a better performance when executing \ac{ML} inference 
tasks. In particular,
the overall layout and flow of the systems, recognizing the role of each
element that takes part in the inference task, as well as the management and 
orchestration of the network, analysing the strategies for scheduling and 
communication methods.

    
\section{System Network Layouts} \label{sec:se31}

Firstly, we will analyse how the proposed frameworks make use of different system 
layouts to achieve collaborative inference in edge and far-edge devices. We will
classify them into two main categories, according to their system layout:
Model Bi-Partitioning and Model Local Partitioning.

% \subsection{Multi-Model Partitioning} \label{sec:se311}

% Rather than focusing on partitioning a single monolithic model, these
% frameworks deploy differently sized \ac{ML} models throughout several
% hierarchies of devices according to their computational capabilities (i.e.
% deploying smaller models closer to the edge and bigger and more capable
% ones closer to the cloud).

% \begin{description}[style=unboxed, leftmargin=1cm, font=\bfseries]

%     \item[N/A] No articles were reviewed yet in this category.

% \end{description}


\subsection{Model Bi-Partitioning} \label{sec:se312}

The approaches reviewed in this category divide the work load in two, deploying 
the initial portion of the inference task to devices closer to the end-user and 
the remaining portion to one or more devices on a higher hierarchical layer, 
further away from the end-user device.

\begin{description}[style=unboxed, leftmargin=1cm, font=\bfseries]

    \item [CoFANN \cite{tran-van_cofann_2025}]
    addresses the trade-off between model accuracy and efficiency by employing
    dynamic inference strategies. CoFANN significantly reduces overall latency
    through search algorithms that minimize additional latency overhead. The 
    inference workload is divided between \ac{IoT} devices and edge devices.
    The network devices are benchmarked in an initial phase, allowing the 
    framework to estimate the time needed for each device on each step of the 
    model. The framework then employs a differentiable search algorithm 
    to find optimal inference strategies for allocating \ac{DNN} tasks to the 
    computing devices. Even though this framework deploys successfully a
    partitioned model between an end device and an edge server, it is limited
    to two-device networks.  

    \item[DINA \cite{t_mohammed_distributed_2020}]
    is a framework that relies on fog nodes to partition inference tasks into
    sub-tasks and deploys them according to each node's characteristics, using
    matching theory for joint partitioning and task offloading. This framework 
    allows for a more efficient use of the network's resources, while also 
    reducing the total computation time.
    This work can process \ac{DNN} models with a \ac{DAG} topology, composed of 
    fully connected and convolutional layers. This framework determines how to 
    partition the inference task and then divides it according to the network's 
    condition. It should be stated that the network is expected to be homogeneous.

    \item [\textcite{ma_joint_2025}]
    proposes a framework that focuses on minimizing energy consumption, making use
    of multiple early-exits along the inference process, when confidence thresholds 
    are met. The framework partitions each model instance into two segments, depending
    on the end device's capabilities. The initial segment is deployed on the end 
    device, while the remaining segment is deployed on an edge server. The framework 
    deploys the same model architecture across a range of devices participating in
    the inference task, adapting each instance to the device's capabilities.
    The early-exit mechanism allows for task termination prematurely in 
    early stages of the inference task, reducing energy consumption and
    latency, given the model is confident enough. 

\end{description}


\subsection{Model Local Partitioning} \label{sec:se313}

In this subsection, we will analyse approaches that deploy the whole partitioned
system to devices locally, without requiring computing power from nodes closer
to the cloud.

\begin{description}[style=unboxed, leftmargin=1cm, font=\bfseries]

    \item[CAS \cite{wang_context-aware_2021}]
    is a framework that successfully adapts a given model in real-time, upon a
    network context change. It generates a "partition state graph", containing 
    the set of possible partitions, given a
    \ac{DNN} model, choosing one based on the current network status, such as 
    inference latency and resources budgets. This work gathers the information
    of the network, taking note on possibly constrained devices. This framework
    creates a fast adapting system, responding to network fluctuations, keeping 
    the active deployment optimized.

    % \item[NAMP \cite{rego::namp}]
    % approaches %TODO

    \item[\textcite{p_dai_joint_2025}] 
    presents a solution to the collaborative inference problem that optimizes
    device placement and model partitioning, accounting for particularities of 
    heterogeneous edge resources and \ac{CNN} models, with the objective to
    maximize throughput.
    A combination of intra-layer and inter-layer partitioning techniques are 
    applied in the proposed approach. This approach reformulates the network as 
    a pipeline system, where each device is responsible for simultaneous 
    inference batches. The \ac{DNN} model segments are tailored for the devices 
    in the network and assigned according to the network topology, that can be 
    represented as a \ac{DAG}. The devices in the network can have sequential 
    or parallel relationships, that affect the assigned model segment.

    \item[\textcite{w_miao_adaptive_2020}]
    explores an approach based on splitting and deploying a given \ac{DNN} model 
    with a \ac{DAG} topology to several devices, exploiting parallelizable 
    segments (branches), based on model data and computation time estimations 
    for each layer and branch processing.
    The characteristics of this network are also taken into account, such as, 
    its bandwidth and the number of devices present in the network. 
    The authors also allow the network to be based on devices with homogeneous 
    computation times or heterogeneous, given the computation time estimations
    of each layer and branch for each device.
    This work also points out the unreliability and unsafety in partitioning the 
    model between data centers and edge devices.

\end{description}

%TODO acabar conclusão de system layouts (incluir ref à tabela)
To better standardize the following synthesis, we will recognize previously 
referred user, end devices and similarly constrained devices as our definition 
of far-edge devices. Likewise, the intermediate devices, placed higher in the
hierarchy than the far-edge devices (e.g.: edge nodes, local servers), will be 
recognized simply as edge devices.

\begin{table}[ht]
\centering
\begin{tabularx}{\textwidth}{|l|Y|Y|Y|}
\hline
\textbf{Strategy} & \textbf{System Layout} & \textbf{System Agents} & \textbf{Partitioning Method} \\
\hline
CoFANN \cite{tran-van_cofann_2025} & Model Bi-Partitioning & Far-Edge Devices, Edge Devices & Graph Partitioning \\
\hline
DINA \cite{t_mohammed_distributed_2020} & Model Bi-Partitioning & Far-Edge Devices, Edge Devices & Graph Partitioning \\
\hline
CAS \cite{wang_context-aware_2021} & Model Local Partitioning & Edge Devices & Graph Partitioning \\
\hline
\textcite{ma_joint_2025} & Model Bi-Partitioning & Far-Edge Devices, Edge Devices & Inter-Layer Partitioning \\
\hline
\textcite{p_dai_joint_2025} & Model Local Partitioning & Edge Devices & Inter and Intra-Layer Partitioning\\
\hline
\textcite{w_miao_adaptive_2020} & Model Local Partitioning & Far-Edge Devices & Graph Partitioning \\
\hline
\end{tabularx}
\vspace{2mm}
\caption{Overview of the system layouts of the reviewed strategies.}\label{tab:sota_layout}
\end{table}


\section{Network Node Management} \label{sec:se32}

In a distributed system, the orchestration and management of every node of a
given network is of utmost importance. Since we are reviewing collaborative
frameworks, we will be analysing how the presented frameworks strategize the 
scheduling and communications between their devices. Furthermore, the frameworks 
will be separated into two classes, heuristical and reinforcement learning 
scheduling approaches.

\subsection{Heuristical Approach} \label{sec:se321}

Predicting the assignments that lead to the most efficient processing of any
given inference task has been revealed to be a highly complex problem. Therefore,
several works have relied on heuristical approaches to simplify the problem
avoiding the loss of a big part of the desired efficiency. In this subsection
we will be reviewing the frameworks that rely on heuristics, which heuristics are
chosen and how the network is managed.

\begin{description}[style=unboxed, leftmargin=1cm, font=\bfseries]

    \item[CAS \cite{wang_context-aware_2021}]
    makes use of "the neighbour effect" and assigns the initial partitions 
    to the first devices, closer to the user, and allocates the tail layers to
    devices closer to the cloud, noted to be heavier on computation time
    even if lighter on the data. All the devices' constraints
    are gathered in a previous stage and taken into account. Network conditions
    are monitored in real-time, allowing the framework to strive for an optimal
    partitioning and deployment, adapting to network changes, such as device power 
    and network bandwidth fluctuations, as well as user requirement of latency. 

    \item [CoFANN \cite{tran-van_cofann_2025}]
    is a minimalist framework, in regards to network management, taking only
    into account a pair of devices. Nonetheless, it looks to minimize inference
    time and data transfers by searching thoroughly a continuous search space
    created from
    factors in the inference process. To minimize the data transfer, the initial
    raw data is processed in the end device and condensed as an intermediate 
    feature tensor to be wirelessly transmitted to an edge server.
    %TODO rever separação layout/flow vs network management
    % FR não comentou


    \item[DINA \cite{t_mohammed_distributed_2020}]
    approaches the scheduling problem using matching games, based on game theory.
    Matching is obtained according to preference relations, similarly to 
    partitioning, which are derived by means of utility functions, representing 
    the benefit obtained by executing a given task on that device. The initial 
    matching is further improved by proposing swapping edge nodes between two 
    user nodes, if it benefits both parties.

    \item[\textcite{p_dai_joint_2025}]
    relies on device placement to assign the model partitions to the devices.
    The feature extractor segment of the \ac{CNN} models is distributed along the
    network, while the classifier, as it requires less computation, is assigned to
    a singular device, given it has sufficient available memory for the corresponding
    weights. Intermediate inference results are exchanged between devices
    through a wireless link, using a "OFDMA\footnote{Orthogonal Frequency-Division 
    Multiple Access}-based wireless transmission model",
    allowing for a higher throughput.
    This solution assumes that the pipeline supports asynchronous communication.

    \item[\textcite{w_miao_adaptive_2020}]
    minimizes the inference time relying on estimations of the computing time
    needed at each step. Since the inference time is determined by the
    edge device that gets results last, they approach scheduling by reducing 
    it to a load balancing problem. The data offloading between devices is 
    implemented by \ac{RPC}s.

\end{description}

\subsection{Reinforcement Learning Approach} \label{sec:se322}

Rather than relying on simple heuristics to achieve a near-optimal
solution, the works in this subsection build reinforcement learning models to 
capture the complex relationship between the network and the optimal assignment
to a greater extent. 

\begin{description}[style=unboxed, leftmargin=1cm, font=\bfseries]

    \item[\textcite{ma_joint_2025}]
    makes use of a \ac{DRL}-based algorithm, employing an actor-critic model to
    determine the optimal resource assignment and inference task's layer to mark
    the bi-partitioning point. The \ac{PPO} algorithm is used to train the model,
    optimizing cost evolution, balancing energy consumption and inference latency.
    This approach takes into account all relevant system information, such as task
    characteristics, device capabilities and network conditions, to make informed
    decisions on resource allocation and task partitioning.

\end{description}

%TODO conclusões de network management (c/ tabela)

\begin{table}[ht]
\centering
\begin{tabularx}{\textwidth}{|l|Y|Y|Y|Y|}
\hline
\textbf{Strategy} & \textbf{Scheduling Approach} & \textbf{Optimized Metrics} & \textbf{Communi-cation Method} & \textbf{Tolerates Node Hete-rogeneity} \\
\hline
CAS \cite{wang_context-aware_2021} & Heuristical & Inference Time & Wireless link (3G, 4G) & \ding{51} \\
\hline
CoFANN \cite{tran-van_cofann_2025} & Heuristical & Inference Time, Data Transferred & Wireless link & \ding{51} \\
\hline
DINA \cite{t_mohammed_distributed_2020} & Heuristical & Inference Time, Resource Utilization & Wireless link & \ding{55} \\
\hline
\textcite{ma_joint_2025} & Reinforcement Learning & Inference Time, Energy Consumption & N/A & \ding{51} \\
\hline
\textcite{p_dai_joint_2025} & Heuristical & Device Placement & OFDMA-based links & \ding{51} \\
\hline
\textcite{w_miao_adaptive_2020} & Heuristical & Inference Time & \ac{RPC} & \ding{51}\\
\hline
\end{tabularx}
\vspace{2mm}
\caption{Overview of the network management in the reviewed strategies.}\label{tab:sota_scheduling}
\end{table}

\section{Final Remarks}

%TODO conclusões gerais do SOTA 
%TODO incluir problema/limitação de cada framework (na tabela?)

\begin{table}[ht]
\centering
\begin{tabularx}{\textwidth}{|p{2.5cm}|X|X|}
\hline
\textbf{Strategy} & \multicolumn{1}{c|}{\textbf{Overview}} & \multicolumn{1}{c|}{\textbf{Relevant Aspects}} \\
\hline
CAS \cite{wang_context-aware_2021} 
& 
    Edge system responsive to network changes, allows for local collaborative
    inference with low latency withstanding network fluctuations. Deployed on
    edge devices, avoiding reliance on large computing resources.
&  
    \begin{itemize}[noitemsep, topsep=0pt, leftmargin=*]
        \item Deployed partition responsiveness to network changes in real-time.
    \end{itemize}
\\
\hline
CoFANN \cite{tran-van_cofann_2025} 
& 
    Proposes a bi-partitioned system that uses end devices in the network to compute 
    initial inference. The framework relies on only two devices, one on-site and
    one edge server. 
&  
    \begin{itemize}[noitemsep, topsep=0pt, leftmargin=*]
        \item Employment of a gradient descent based search algorithm to detect the optimal
        inference strategy, to minimize latency and data transfer.  
    \end{itemize}
\\
\hline
DINA \cite{t_mohammed_distributed_2020} 
& 
    Employs a bi-hierarchical fog computing architecture to execute inference tasks.
    Deploys partitioned \ac{DNN} models with \ac{DAG} topologies to homogeneous
    devices. Adapts to reduce execution time in varying network conditions.
&  
    \begin{itemize}[noitemsep, topsep=0pt, leftmargin=*]
        \item Extension of a local inference task to a distributed edge device network.
        \item Partition and compression techniques for \ac{CNN} models.
        \item Matching theory for scheduling purposes.
    \end{itemize}
\\
\hline
\textcite{ma_joint_2025} 
& 
    Framework that minimizes energy consumption in inference networks by employing
    early-exit mechanisms and bi-partitioning of \ac{DNN} models. The framework
    relies on a \ac{DRL}-based algorithm to optimize resource assignment and
    partitioning point.
&  
    \begin{itemize}[noitemsep, topsep=0pt, leftmargin=*]
        \item Use of early-exit mechanisms to reduce energy consumption and latency.
        \item \ac{DRL}-based scheduling approach using \ac{PPO} algorithm.
    \end{itemize}
\\
\hline
\textcite{p_dai_joint_2025} 
& 
    Widely applicable framework for collaborative inference in edge networks,
    optimizing throughput and adapting partitioning according to each device's 
    capabilities. The framework has been shown to be a highly scalable,
    robust and efficient solution.
& 
    \begin{itemize}[noitemsep, topsep=0pt, leftmargin=*]
        \item Partition and parallelization techniques for \ac{CNN} models.
        \item Device placement and network topology considerations for scheduling. 
    \end{itemize}
\\
\hline
\textcite{w_miao_adaptive_2020} 
& 
    Collaborative inference framework that takes advantage of parallelizable
    segments in \ac{DNN} models with \ac{DAG} topologies. The framework
    adapts to heterogeneous devices, balancing the load between them to
    minimize the overall inference time.
& 
    \begin{itemize}[noitemsep, topsep=0pt, leftmargin=*]
        \item Load balancing approach for heterogeneous devices.
    \end{itemize}
\\
\hline
\end{tabularx}
\vspace{2mm}
\caption{Overview of the reviewed strategies.}\label{tab:sota_overview}
\end{table}
